\chapter*{前言}
\addcontentsline{toc}{chapter}{前言}
\markboth{前言}{前言}


\section*{為什麼要邀請你讀這本書？}
\normalsize
也許你從來沒有聽過啟示錄，或許你知道啟示錄是聖經的最後一本書卷，可是懷疑它的真實性。也許你常常讀聖經，只是對啟示錄比較陌生。
如果你正有這樣的疑問，那這本書是專門爲你而寫的。

如果你未接觸基督信仰，你可能會以為啟示錄只是給基督徒讀的。可是我想告訴你，啟示錄是神自己邀請你來認識祂的福音的書。約翰寫下啟示錄之後，啟示錄就在小亞細亞眾教會中流傳開來。許多基督徒在壓力與逼迫之下，仍然持守信仰，甚至為基督作見證直到犧牲性命。這些早期基督徒所相信的，是一位真實來到世上的基督。

\section*{寫作動機—為什麼要寫這本書？}

因為關於末世預言的電影很多，也有不少電影和動畫以末世為背景，使人對啟示錄所描繪的世界產生好奇。因此，有些人也會翻閱啟示錄，想看看聖經究竟有沒有記載「世界末日會發生什麼事」。曾經有一位青少年告訴我，他讀聖經的目的，就是想看看裡面到底有沒有記載世界末日會發生什麼事。
 
當基督徒閱讀啟示錄時，我們的焦點也可能落在如何解釋龍、獸、印記、千禧年等問題上。其實，這也是我自己閱讀啟示錄時所經歷的一個轉折。因此，我寫這本書，是希望與讀者重新走進啟示錄。不是先從每一個單一事件開始，而是從耶穌基督開始。啟示錄的第一句話，不就是「耶穌基督的啟示」嗎？它是耶穌自己對你我說話。它讓我們知道，祂掌管歷史，祂審判邪惡，祂勝過撒但與死亡，祂再來，祂更新萬有，最後把祂所救贖的百姓帶入永恆的榮耀。

同時，我也盼望這本書能成為協助主日學教師的材料。因此，我盡量交代經文的背景與信息。盼望無論是個人靈修、小組查經，或教師教學，都能藉此一同思想啟示錄所見證的耶穌基督。

\section*{基督／羔羊是中心—如何讀啟示錄？}

當你翻閱啟示錄時，可能會很困惑，耶穌不是人嗎？怎麼就變成羔羊了？
這正是我們需要認識聖經整體的地方。耶穌第一次來，是贖罪的羔羊，祂被人類想出來最殘酷的刑罰，被釘在十字架上。這是他主動的行為，為了贖回在罪中被綑綁的人。我們哪裡知道祂是為我們的過犯而受害呢？聖經的《以賽亞書》（約是耶穌來到地上前700年），早已預言祂的代死。祂為我們的罪孽壓傷，「因祂受的刑罰，我們得平安；因祂受的鞭傷，我們得醫治。」（以賽亞書53:5）原來祂像是一隻羊羔被獻祭。

然而故事就此結束了嗎？沒有。啟示錄接續了耶穌的故事。祂以羔羊的身份來到，但如今是站立得勝的羔羊，並且是領軍的君王，我們也可以認得祂，因為祂穿著濺了血的衣服。（啟示錄19:13）所以啟示錄是一個以基督為中心的故事。祂掌權，展開書卷，審判邪惡，擊敗撒但，並且將祂所救贖的百姓帶入新創造的榮耀。

歷世歷代數不清的基督徒何其有幸，能被恩典帶入神的榮耀裡。相信正在讀本書的你也是其中大蒙恩典的一位，並且希望你看完此書之後，願意接受主恩典的邀請。

\section*{書名：「被恩典帶入榮耀」}

此書承襲《被恩典追上》的命名：人先被神的恩典追上，繼而被神的恩典改變，最後被神的恩典帶入祂的榮耀。因此，這一系列的名稱回應了一個蒙恩得救的基督徒所走的道路：
\begin{itemize}
\item 「被恩典追上」
\item 「被恩典改變」
\item 「被恩典帶入榮耀」
\end{itemize}
此書封面的「恩典」二字是紅色，象徵耶穌基督的寶血帶我們進入祂的榮耀裡。
那麼，為什麼要先寫第三部《被恩典帶入榮耀》呢？因為\textbf{知道終點，才能更清楚地知道改變的目的。}神在基督裡的救恩有一個榮耀的終點，就是使我們最終與基督同得榮耀。
當我們知道自己將往哪裡去，就更能明白今天為什麼需要被神改變。這正是保羅在《羅馬書》8章所描繪的救恩金鏈：

\begin{quote}
「因為祂\textbf{預先所知道}的人，就『預定』效法祂兒子的模樣，使祂兒子在許多弟兄中作長子。預先所定下的人又『召』他們來；所召來的人又『稱他們為義』；所稱為義的人又叫他們『得榮耀』。」—羅馬書 8:29-30
\end{quote}
從預定、呼召、稱義，到最終的得榮耀，這條救恩的道路不是人自己走向榮耀，而是神以祂的恩典一路帶領祂所救贖的人，直到榮耀的終點。

\section*{研究脈絡與主要著作}

歷來對啟示錄的研究卷帙浩繁，我在準備這本導讀性質的書時，只能有所選擇。
本書除了參考多位解經家與學者的研究，也整理了我過去研讀啟示錄以及教導主日學時所累積的資料與心得。
並在這次重新整理的過程中，我有許多新的發現及收穫，使我對啟示錄的基督，有了更深一層的認識。

以下幾位作者的書是我寫此書的參考重點：
\begin{enumerate}
\item John MacArthur（約翰. 麥克阿瑟）{\footnotesize\footnote{MacArthur, J. (2007). \textit{Because the Time Is Near: John MacArthur Explains the Book of Revelation}. Moody Publishers.}} —他的逐段解經，以及他對基督再來、基督掌權與最終得勝的論述。

\item Grant R. Osborne（格蘭特. 奧斯本）\footnote{Osborne, G. R. (2002). \textit{Revelation}. Baker Academic.} —他對不同啟示錄解讀的探討，幫助我從不同角度思考本書的象徵、結構與末世論主題。
    
\item G. K. Beale（G. K.\ 比爾）\footnote{Beale, G. K. (1999). \textit{The Book of Revelation: A Commentary on the Greek Text}. Wm. B. Eerdmans Publishing Co.} —他對舊約背景的廣泛研究，幫助我將啟示錄放在整體聖經故事中來理解，並看見其中豐富的舊約意象與神學脈絡。
    
\item Craig R. Koester（克雷格. 科斯特）\footnote{Koester, C. R. (2014). \textit{Revelation: A New Translation with Introduction and Commentary}. Yale University Press.} —他對啟示錄歷史背景與文學脈絡，使我更能理解這卷書向最初眾教會所傳達的信息。
    
\item Richard Bauckham（理查. 鮑坎）\footnote{Bauckham, R. (1993). \textit{The Theology of the Book of Revelation}. Cambridge University Press.} 
—他從神學角度的研究，幫助我進一步思考許多核心主題。
\end{enumerate}

這些作者各有不同的貢獻以及詮釋。因此，我盡力地呈現主要的不同觀點。本書的目的，是以容易理解的方式帶領讀者認識啟示錄，同時始終把它的核心聚焦在耶穌基督的得勝，以及祂的百姓最終進入永恆的榮耀。\footnote{
我也感謝吳榮滁牧師所開發的「新約引用舊約分析平台」，彷彿是神學院的教授在我旁指導，並使我查考的時間，事半功倍。https://chatgpt.com/g/g-69e26d759c108191ba873085c5722e4c-xin-yue-yin-yong-jiu-yue-fen-xi-ping-tai-nt-ot-citation-analyzer

我並感謝蘇桂村牧師二十多年來亦師亦友的陪伴與代禱。在他各地講道百忙之中，他仍願意逐頁讀我的寫作，提出許多寶貴的意見與建議，使我獲益良多。對於他的耐心與付出，我深表感謝。

當然，本書若有任何內容上的錯誤或疏漏，仍由我本人負全部責任。
}

\section*{表格與時間線：協助快速閱讀}

啟示錄包含豐富的異象與神學主題。為了幫助現代讀者理解並掌握這些豐富的內容，凡是可以清楚地透過表格或時間線呈現的資訊，我都盡量以這樣的方式加以整理。

如此一來，正文的討論就可以持續聚焦於耶穌基督的掌權、得勝，以及祂最終更新萬有的工作。對於經文中不同的解釋觀點，我在註腳或是附錄中加以說明。

\section*{本書的插圖：進入啟示錄的異象}

啟示錄充滿豐富而鮮明的圖像：龍、獸、羔羊、金燈臺、號筒、印、碗，以及新耶路撒冷等。對現代讀者而言，這些異象若只透過文字描述，有時不容易形成具體的畫面。因此，本書適度加入相關插圖，幫助讀者更直觀地掌握經文中的場景與信息。

雖然插圖的構思與製作需要投入不少時間，但若能幫助弟兄姊妹更深入地進入經文所呈現的異象，這些努力便是值得的。

然而，插圖也可能影響讀者對異象的想像，甚至加入經文未明確提供的細節。為此，本書的插圖盡量保持簡潔，主要發揮視覺輔助的作用，而不取代經文本身的描述與解釋。\footnote{本書的插圖由作者用 Google Gemini 反覆修改製作而定稿，並經過作者進一步的審閱、編輯與篩選。}

讀者應當以啟示錄的經文為依據，將插圖視為理解人物、場景與信息的輔助工具，而非解釋經文的權威。


\section*{結語}

寫完初稿時，我的心情是敬畏的。我這個蒙恩得救的人，竟要盡力寫出祂的榮耀。祂慈愛的大光照著我，使我的筆快活如馬，廢寢忘餐；可是，我這支不達意的筆，在我的造物主面前，又笨拙不如烏鴉一聲驚叫。

然而，我的造物主是知道的。祂知道我筆下所要述說的，是那位我所愛、所敬拜、所仰望的主。正如使徒約翰在啟示錄最後所說的：

\begin{quote}
「主耶穌啊，我願你來！」
\end{quote}

\vspace{2em}

\begin{flushright}
主僕  \\
郭景光弟兄\\
紐約新城歸正教會長老  \\
紐約州艾姆赫斯特
\end{flushright}